% GENERATED FILE. Edit canonical lessons, then run npm run generate:books.
% canonical-source-hash: fnv1a64:f0cf0839838b9192
% canonical-lessons: IT-C487-masticare, IT-C487-digerire, IT-C487-malessere, IT-C487-vomitare, IT-C487-vertigini

\chapter{Stomaco e testa}
\label{ch:it-stomaco-e-testa}
\hlchaptermodality{487}

\begin{chapteropening}
\textbf{By the end of this chapter:} \emph{I can describe stomach trouble and dizziness.}

Complete the last lesson of chapter 487: I can describe stomach trouble and dizziness.
\end{chapteropening}

\section[masticare]{masticare — to chew}

\label{lesson:IT-C487-masticare}



\begin{quote}
Before the new one: say the Italian for allergy, then the Italian for sting.
\end{quote}

\subsection*{You'll want to know: masticare}
\textbf{masticare} — \textquotedblleft{}to chew\textquotedblright{}.

\begin{quote}
Mastica bene prima di mandare giù.
\end{quote}

\begin{quote}
\emph{Chew well before you swallow.}
\end{quote}

\emph{La gomma da masticare} — chewing gum.

\begin{cousinweb}
\emph{masticare} is inherited from Late Latin \emph{masticāre} \textquotedblleft{}to chew\textquotedblright{} (EtymDB; Lewis \& Short), usually linked to Greek \emph{mastikhân} \textquotedblleft{}to gnash the teeth\textquotedblright{}.
\end{cousinweb}

\begin{grammarlens}[title={What you've built}]
One more word for telling, explaining and describing at B1.
\end{grammarlens}

\subsection*{Your turn}
\begin{itemize}
\raggedright
  \item \emph{Say it:} \emph{masticare}
  \item \emph{Say it:} \emph{masticare}, once more
  \item \emph{Say it:} a sentence of your own with \emph{masticare}
  \item \emph{Recall:} say the Italian for allergy, then the Italian for sting, then say \emph{masticare} again
\end{itemize}

\begin{tcolorbox}[breakable,colback=teal!4,colframe=teal!35!black,title={Before you move on}]
What does masticare mean? (\textquotedblleft{}To chew\textquotedblright{}.) Say it once more.
\end{tcolorbox}
\section[digerire]{digerire — to digest}

\label{lesson:IT-C487-digerire}



\begin{quote}
Before the new one: say the Italian for sting, then the Italian for to chew.
\end{quote}

\subsection*{You'll want to know: digerire}
\textbf{digerire} — \textquotedblleft{}to digest\textquotedblright{}.

\begin{quote}
Non ho digerito bene la cena.
\end{quote}

\begin{quote}
\emph{I did not digest dinner well.}
\end{quote}

Also figurative: \emph{non lo digerisco} — I cannot stand him. An \emph{-isc-} verb.

\begin{cousinweb}
\emph{digerire} is a learned borrowing from Latin \emph{dīgerere} \textquotedblleft{}to divide, to dissolve\textquotedblright{} (EtymDB), from \emph{dis-} and \emph{gerere} \textquotedblleft{}to carry\textquotedblright{} (Lewis \& Short).
\end{cousinweb}

\begin{grammarlens}[title={What you've built}]
One more word for telling, explaining and describing at B1.
\end{grammarlens}

\subsection*{Your turn}
\begin{itemize}
\raggedright
  \item \emph{Say it:} \emph{digerire}
  \item \emph{Say it:} \emph{digerire}, once more
  \item \emph{Say it:} a sentence of your own with \emph{digerire}
  \item \emph{Recall:} say the Italian for sting, then the Italian for to chew, then say \emph{digerire} again
\end{itemize}

\begin{tcolorbox}[breakable,colback=teal!4,colframe=teal!35!black,title={Before you move on}]
What does digerire mean? (\textquotedblleft{}To digest\textquotedblright{}.) Say it once more.
\end{tcolorbox}
\section[il malessere]{il malessere — feeling unwell; unease}

\label{lesson:IT-C487-malessere}



\begin{quote}
Before the new one: say the Italian for to chew, then the Italian for to digest.
\end{quote}

\subsection*{You'll want to know: il malessere}
\textbf{il malessere} — \textquotedblleft{}feeling unwell; unease\textquotedblright{}.

\begin{quote}
Dopo pranzo ho avuto un leggero malessere.
\end{quote}

\begin{quote}
\emph{After lunch I felt a little unwell.}
\end{quote}

A vague feeling that something is wrong, in the body or the mind.

\begin{cousinweb}
\emph{malessere} is built from \emph{male} \textquotedblleft{}badly\textquotedblright{}, from Latin \emph{male}, and \emph{essere} \textquotedblleft{}to be\textquotedblright{}, from Latin \emph{esse} (Lewis \& Short): a being-unwell.
\end{cousinweb}

\begin{grammarlens}[title={What you've built}]
One more word for telling, explaining and describing at B1.
\end{grammarlens}

\subsection*{Your turn}
\begin{itemize}
\raggedright
  \item \emph{Say it:} \emph{il malessere}
  \item \emph{Say it:} \emph{il malessere}, once more
  \item \emph{Say it:} a sentence of your own with \emph{il malessere}
  \item \emph{Recall:} say the Italian for to chew, then the Italian for to digest, then say \emph{il malessere} again
\end{itemize}

\begin{tcolorbox}[breakable,colback=teal!4,colframe=teal!35!black,title={Before you move on}]
What does il malessere mean? (\textquotedblleft{}Feeling unwell; unease\textquotedblright{}.) Say it once more.
\end{tcolorbox}
\section[vomitare]{vomitare — to vomit, to be sick}

\label{lesson:IT-C487-vomitare}



\begin{quote}
Before the new one: say the Italian for to digest, then the Italian for feeling unwell.
\end{quote}

\subsection*{You'll want to know: vomitare}
\textbf{vomitare} — \textquotedblleft{}to vomit, to be sick\textquotedblright{}.

\begin{quote}
Ho mangiato troppo e ho vomitato.
\end{quote}

\begin{quote}
\emph{I ate too much and was sick.}
\end{quote}

More politely: \emph{rimettere}.

\begin{cousinweb}
\emph{vomitare} is a learned borrowing from Latin \emph{vomitāre} \textquotedblleft{}to vomit often\textquotedblright{} (EtymDB), from \emph{vomere} \textquotedblleft{}to vomit\textquotedblright{} (Lewis \& Short).
\end{cousinweb}

\begin{grammarlens}[title={What you've built}]
One more word for telling, explaining and describing at B1.
\end{grammarlens}

\subsection*{Your turn}
\begin{itemize}
\raggedright
  \item \emph{Say it:} \emph{vomitare}
  \item \emph{Say it:} \emph{vomitare}, once more
  \item \emph{Say it:} a sentence of your own with \emph{vomitare}
  \item \emph{Recall:} say the Italian for to digest, then the Italian for feeling unwell, then say \emph{vomitare} again
\end{itemize}

\begin{tcolorbox}[breakable,colback=teal!4,colframe=teal!35!black,title={Before you move on}]
What does vomitare mean? (\textquotedblleft{}To vomit, to be sick\textquotedblright{}.) Say it once more.
\end{tcolorbox}
\section[le vertigini]{le vertigini — dizziness, vertigo}

\label{lesson:IT-C487-vertigini}



\begin{quote}
Before the new one: say the Italian for feeling unwell, then the Italian for to vomit.
\end{quote}

\subsection*{You'll want to know: le vertigini}
\textbf{le vertigini} — \textquotedblleft{}dizziness, vertigo\textquotedblright{}.

\begin{quote}
Quando guardo giù dal balcone ho le vertigini.
\end{quote}

\begin{quote}
\emph{When I look down from the balcony I feel dizzy.}
\end{quote}

\emph{Avere le vertigini} — to feel dizzy; to be afraid of heights is \emph{soffrire di vertigini}.

\begin{cousinweb}
\emph{vertigini} is the plural of \emph{vertigine}, a learned borrowing from Latin \emph{vertīgō} \textquotedblleft{}a turning round, dizziness\textquotedblright{} (Lewis \& Short), from \emph{vertere} \textquotedblleft{}to turn\textquotedblright{}.
\end{cousinweb}

\begin{grammarlens}[title={What you've built}]
One more word for telling, explaining and describing at B1.
\end{grammarlens}

\subsection*{Your turn}
\begin{itemize}
\raggedright
  \item \emph{Say it:} \emph{le vertigini}
  \item \emph{Say it:} \emph{le vertigini}, once more
  \item \emph{Say it:} a sentence of your own with \emph{le vertigini}
  \item \emph{Recall:} say the Italian for feeling unwell, then the Italian for to vomit, then say \emph{le vertigini} again
\end{itemize}

\begin{tcolorbox}[breakable,colback=teal!4,colframe=teal!35!black,title={Before you move on}]
What does le vertigini mean? (\textquotedblleft{}Dizziness, vertigo\textquotedblright{}.) Say it once more.
\end{tcolorbox}
